\chapter{Pré-requisitos}

Este capítulo descreve os conhecimentos prévios recomendados e o ambiente necessário
para acompanhar o minicurso e reproduzir os exemplos práticos de monitoramento de
funções virtualizadas de rede com eBPF.

% ----------------------------------------------------------------
\section{Conhecimentos necessários}
% ----------------------------------------------------------------

O minicurso foi concebido para participantes com familiaridade básica com sistemas
operacionais Linux, programação em linguagem C e fundamentos de redes de computadores.
Não é exigida experiência prévia com desenvolvimento de kernel, com eBPF ou com
funções virtualizadas de rede (NFV).
Os tópicos a seguir representam a base mínima esperada.

\textbf{Linha de comando Linux.}
O participante deve ser capaz de navegar pelo sistema de arquivos, executar comandos
com privilégio de superusuário (\texttt{sudo}), instalar pacotes com o gerenciador
de pacotes da distribuição e redirecionar a saída de programas para arquivos e
para outros comandos via \textit{pipe} \cite{manpages2024}.

\textbf{Linguagem C --- nível introdutório.}
Programas eBPF são escritos em um subconjunto restrito de C e compilados para
\textit{bytecode} BPF \cite{linuxkernel2024}.
É suficiente compreender tipos primitivos, ponteiros, estruturas (\texttt{struct})
e a distinção entre espaço de usuário e espaço de kernel.
Conceitos avançados, como alocação dinâmica de memória ou programação multithread,
não são necessários neste nível introdutório.

\textbf{Fundamentos de redes TCP/IP.}
Para os exemplos de monitoramento de rede, é necessário compreender o modelo de
camadas TCP/IP, o conceito de pacote, como endereços IP e portas identificam
conexões e a diferença entre protocolos orientados a conexão (TCP) e não orientados
(UDP) \cite{gregg2019}.
Esse conhecimento contextualiza os pontos de observação que programas eBPF ocupam
na pilha de rede do kernel.

\textbf{Conceitos básicos de sistemas operacionais.}
É recomendável ter noção do que é uma chamada de sistema (\textit{syscall}), como
processos são escalonados e qual o papel do kernel na mediação entre aplicações e
o hardware \cite{gregg2020}.
Em particular, compreender o que é um \textit{file descriptor} e como o sistema
de arquivos virtual \texttt{/proc} expõe informações de processos é útil para
interpretar os dados coletados pelos programas eBPF \cite{manpages2024}.

\textbf{Noções de virtualização de funções de rede (NFV).}
Os exemplos práticos são contextualizados no cenário de monitoramento de funções
virtualizadas de rede conforme o arcabouço do ETSI NFV \cite{etsi2014}.
Não é necessário ter experiência prévia com NFV, mas compreender que uma função
de rede virtualizada é um software executando sobre infraestrutura de propósito
geral --- em vez de hardware dedicado --- facilita a interpretação dos resultados
obtidos \cite{venancio2021}.

% ----------------------------------------------------------------
\section{Ambiente Linux}
% ----------------------------------------------------------------

Todos os exemplos foram desenvolvidos e testados em ambiente Linux com kernel
versão 5.15 ou superior, que é o mínimo recomendado para suporte estável às
funcionalidades de eBPF utilizadas no minicurso, incluindo
\texttt{BPF\_MAP\_TYPE\_RINGBUF} e suporte completo a BPF CO-RE
\cite{linuxkernel2024}.
O código-fonte e os scripts de configuração do ambiente estão disponíveis no
repositório público do minicurso \cite{plinta2026repo}.

\textbf{Distribuição recomendada.}
Recomenda-se Ubuntu 22.04 LTS ou Debian 12, por oferecerem pacotes atualizados
de ferramentas eBPF nos repositórios oficiais e ampla documentação de suporte.
Outras distribuições baseadas em Debian ou em Red Hat (Fedora, RHEL 9+) são
igualmente compatíveis, bastando adaptar os comandos de instalação de pacotes.

\textbf{Uso em máquinas virtuais.}
O eBPF pode ser utilizado em máquinas virtuais, desde que o kernel do sistema
convidado tenha sido compilado com as opções \texttt{CONFIG\_BPF},
\texttt{CONFIG\_BPF\_SYSCALL} e \texttt{CONFIG\_BPF\_JIT} habilitadas
\cite{linuxkernel2024}.
Para fins didáticos, recomenda-se o uso de uma máquina virtual com acesso
irrestrito ao kernel, que oferece maior visibilidade sobre o sistema e evita
restrições impostas por camadas de isolamento adicionais.

\textbf{Verificação da versão do kernel.}
Antes de iniciar as atividades práticas, recomenda-se confirmar a versão do kernel
em execução e verificar se o sistema de arquivos de depuração está montado:

\begin{lstlisting}[language=bash,
  caption={Verificação de versão do kernel e disponibilidade de tracepoints.}]
uname -r
ls /sys/kernel/debug/tracing/events/syscalls | head -5
\end{lstlisting}

Caso o diretório \texttt{/sys/kernel/debug} não exista ou esteja vazio, o sistema
de arquivos de depuração deve ser montado manualmente:

\begin{lstlisting}[language=bash,
  caption={Montagem do sistema de arquivos de depuração.}]
sudo mount -t debugfs debugfs /sys/kernel/debug
\end{lstlisting}

% ----------------------------------------------------------------
\section{Configuração do ambiente eBPF}
% ----------------------------------------------------------------

Esta seção descreve a instalação das dependências necessárias para compilar e
executar os programas eBPF do minicurso.

\textbf{Instalação das dependências.}
Os seguintes pacotes devem estar presentes no sistema antes do início das
atividades práticas:

\begin{itemize}
  \item \texttt{clang} e \texttt{llvm} --- compilador utilizado para gerar
        \textit{bytecode} BPF a partir de código C \cite{linuxkernel2024};
  \item \texttt{libbpf-dev} --- biblioteca de suporte a programas eBPF com
        suporte a CO-RE \cite{libbpf};
  \item \texttt{bpftrace} --- linguagem de rastreamento de alto nível para
        investigação interativa \cite{bpftrace};
  \item \texttt{bpftool} --- utilitário oficial para inspeção de programas,
        mapas e objetos BPF carregados no kernel \cite{linuxkernel2024};
  \item \texttt{linux-headers-\$(uname -r)} --- cabeçalhos do kernel necessários
        para compilação de programas eBPF;
  \item \texttt{make} e \texttt{gcc} --- utilitários de construção (\textit{build}).
\end{itemize}

Em distribuições baseadas em Debian, a instalação pode ser realizada com:

\begin{lstlisting}[language=bash,
  caption={Instalação das dependências em Ubuntu 22.04 / Debian 12.}]
sudo apt update
sudo apt install clang llvm libbpf-dev bpftrace bpftool \
                 linux-headers-$(uname -r) make gcc
\end{lstlisting}

\textbf{Obtenção do repositório do minicurso.}
Os exemplos práticos estão organizados no repositório \cite{plinta2026repo}.
Para cloná-lo e acessar o diretório de exemplos:

\begin{lstlisting}[language=bash,
  caption={Clonagem do repositório e acesso ao diretório de exemplos.}]
git clone https://github.com/joaopedroplinta/ebpf-scitec.git
cd ebpf-scitec
\end{lstlisting}

\textbf{Verificação da instalação.}
Após instalar as dependências, recomenda-se executar um programa eBPF mínimo com
o bpftrace para confirmar que o ambiente está funcional:

\begin{lstlisting}[language=bash,
  caption={Verificação do ambiente com bpftrace.}]
sudo bpftrace -e 'BEGIN { printf("eBPF funcionando!\n"); exit(); }'
\end{lstlisting}

A saída esperada é \texttt{eBPF funcionando!}, seguida do encerramento do processo.
Caso ocorra erro de permissão, verifique se o comando está sendo executado com
\texttt{sudo} ou se o usuário possui a capacidade \texttt{CAP\_BPF}
\cite{linuxkernel2024}.

\textbf{Instalação das ferramentas BCC.}
Para os exemplos que utilizam o BCC \cite{bcc}, o pacote \texttt{bpfcc-tools}
disponibiliza as ferramentas prontas para uso sem necessidade de compilação:

\begin{lstlisting}[language=bash,
  caption={Instalação das ferramentas BCC.}]
sudo apt install bpfcc-tools python3-bpfcc
\end{lstlisting}

Após a instalação, ferramentas como \texttt{execsnoop-bpfcc} e
\texttt{opensnoop-bpfcc} podem ser executadas diretamente, e os scripts de
monitoramento de rede --- como \texttt{tcpconnect-bpfcc} e
\texttt{tcpretrans-bpfcc} --- estão disponíveis para uso imediato nos exemplos
práticos do minicurso.